\subsection{Deep Learning Ablation Study}
\label{sec:dl_ablations}
To isolate the empirical contributions of backbone architecture, fine-tuning depth, image resolution, head regularisation, data augmentation components, and class imbalance mitigation, we conducted a pre-registered one-factor-at-a-time (OFAT) ablation study comprising 13 configurations across the 5 frozen training cross-validation folds (65 evaluated folds, evaluated on Apple Silicon GPU hardware).
The baseline model---a frozen MobileNetV2~\cite{sandler2018mobilenetv2} backbone with a linear classification head, $p=0.2$ dropout, and comprehensive data augmentation at $224 \times 224$ resolution---achieved a mean Macro~F1 of $0.8515 \pm 0.0235$ (87.19\% accuracy).
Under our pre-registered selection rule ($|\Delta| > \sigma = 0.0235$), swapping the frozen backbone to EfficientNet-B0~\cite{tan2019efficientnet} ($0.8557 \pm 0.0263$), adjusting dropout ($p=0.0$: $0.8503$, $p=0.5$: $0.8479$), omitting individual augmentations (e.g., removing $\pm 10^\circ$ rotation: $0.8490$), removing all augmentation ($0.8339 \pm 0.0204$, $\Delta = -0.0176$), and incorporating class weighting ($0.8453$) or balanced sampling ($0.8419$) yielded score differences within $\sigma$, exhibiting no clear difference under the pre-registered rule in overall Macro~F1.
However, class-weighted loss and balanced sampling explicitly improved mean per-fold trash recall from 65.63\% to 72.89\% and 72.84\% respectively (pooled out-of-fold recall from 63/96 to 70/96).
Conversely, reducing input resolution to $160 \times 160$ fell just beyond $\sigma$ (worse, Macro~F1 $0.8271 \pm 0.0257$, $\Delta = -0.0244$).
Crucially, partial fine-tuning of only the final two structural blocks in MobileNetV2 (\texttt{A2.1}, training 893,766 parameters while keeping the first 17 blocks frozen, with differential learning rates of $10^{-3}$ on the head and $10^{-4}$ on the backbone) produced a clear improvement under the rule to \textbf{0.8825 $\pm$ 0.0219 Macro~F1} (89.80\% accuracy, $\Delta = +0.0310 > \sigma$, median best epoch 13), driven by recall gains on trash (79.21\% vs.\ 65.63\%) and metal (94.39\% vs.\ 89.50\%).
We note as a limitation that most frozen configurations reached their best epoch at 38--40 of the pre-registered 40-epoch cap (e.g., \textsc{Base} median best 38, \texttt{A4.2} and \texttt{A5.1} median best 40) whereas \texttt{A2.1} early-stopped naturally (mean 20.2 epochs run); part of \texttt{A2.1}'s margin may reflect this budget constraint, and extended frozen training is left as future work.
Configuration \texttt{A2.1} is on par with the Phase 1 SVM on handcrafted features ($0.8803 \pm 0.0215$, $\Delta = 0.0022 \ll \sigma$), while Daksh's early feature fusion hybrid H2 ($0.9044 \pm 0.0332$) shows no clear gain over \texttt{A2.1} under the $\sigma$ rule ($\Delta = 0.0219 < 0.0332$), with the final evaluation deferred to upcoming late fusion (H3) and the held-out test set.
Consequently, configuration \texttt{A2.1} was designated as the primary deep-learning model, and its out-of-fold predictions were exported for late fusion.

